\subsection{Justification of Algorithmic Choices}

The system developed in Sections~4.1.3 through~4.1.6 represents a deliberate departure from existing approaches. Each algorithm provides the minimal necessary mathematical solution to overcome the three structural barriers identified in Section~4.1.1, which include static variability in patients, the time factor in clinical procedures, and the conversion of numeric outputs to human cognition.

\textbf{K-Prototypes for Population Stratification.} Clustering can be seen as operating in a feature space, which inherently includes both continuous-valued features and discrete categorical features. As discussed in Section~4.1.3, $k$-means algorithm is mathematically ill-defined on categorical attributes, because the average of numerically encoded categories is not meaningful and induces unnecessary geometric structure, whereas $k$-modes loses the notion of variance of continuous attributes by artificially converting them into categorical features. While DBSCAN density-based clustering needs a universal global distance metric and struggles with the "curse of dimensionality" in sparse spaces, Gaussian Mixture Model requires parametric continuous distribution that categorical attributes do not satisfy. Lastly, hierarchical clustering based on Gower's distance, though independent of any distance measure, scales as $O(n^2)$ and lacks any prototype representation, therefore, does not produce any mathematical entity to characterize each cluster. Therefore, the only method among them meeting all of these criteria is the K-Prototypes algorithm, which (i) provides a mathematically well-defined dissimilarity measure using Equation~\eqref{eq:kproto_dist}, (ii) offers a closed-form solution for prototype update being provably optimal w.r.t.\ objective function of each individual variable, and (iii) gives a non-increasing sequence that converges to a local optimum of Equation~\eqref{eq:kproto_obj}.

\textbf{Complementary Validation Criteria.} These six criteria, described in Section~4.1.4, were specifically selected to provide non-redundant, complementary insights about the dataset under investigation. Silhouette Score measures individual data points' geometric fit, whereas Gower's variant allows for measuring geometric fit within the mixed-type space. At the macro level, Calinski-Harabasz index assesses cluster separation based on variance ratios, while Davies-Bouldin Index identifies the most severe instances of overlap that may be obscured by averaging across clusters. Finally, because geometric distances cannot isolate feature-level importance, Cramér's $V$ and $\eta^2$ are employed to measure the statistical effect sizes of individual categorical and numerical variables.

\textbf{PrefixSpan for Trajectory Mining.} For algorithms performing sequential pattern mining, the main contenders are GSP, SPADE, and PrefixSpan. The former is based on the classic generate and test framework akin to Apriori, whereby the growth of the candidate set becomes exponential with respect to the alphabet size and requires multiple scans through the entire dataset at each level of mining. SPADE uses a vertical representation in terms of ID list, which although enabling efficient joining, suffers from severe memory overhead when expanding ID lists for long sequences with large clinical alphabets. PrefixSpan, however, manages to avoid generating candidates altogether with prefix projection technique; the size of the projected dataset $\mathbb{S}|_\alpha$ is strictly non-increasing as the length of the sequence $\alpha$ grows. Together with the inherent data compression due to closed patterns (Equation~\eqref{eq:closed}) and geometric-mean lift re-ranking, which mitigates the effect of frequent item bias, the only possible choice is PrefixSpan algorithm.

